% !TeX program = xelatex
% Jednoducha LaTeX sablona seminární práce podle stylu MGO.
% Kompilace: xelatex seminarni_prace_template.tex (2x kvuli obsahu a poctu stran)

\documentclass[12pt,a4paper]{article}

% --- Jazyk, typografie, fonty ---
\usepackage{fontspec}
\usepackage[czech]{babel}
\usepackage{microtype}
\usepackage{setspace}
\usepackage{ragged2e}
\usepackage{indentfirst}
\usepackage{etoolbox}

% ---------------------------------------------------------
% FONTOVÁ ANALÝZA OFICIÁLNÍ MGO PŘEDLOHY
% Zjištěno přímo z PDF pomocí pdffonts/pdftohtml:
% - titulní strana:
%   škola, autor, předmět, místo/rok .......... Cambria Regular 16 pt, patkové
%   hlavní název .............................. Calibri Regular 28 pt, bezpatkové, NE tučné
%   podnázev ................................... Calibri Regular 16 pt, bezpatkové, NE tučné
% - běžný text a obsah ........................ Cambria Regular 12 pt, patkové
% - nadpis Obsah a nadpisy 1./2./3. úrovně ..... Calibri Bold, bezpatkové, tučné
% - záhlaví .................................... Cambria Italic 12 pt, patkové kurzívní
% - zápatí s číslem stránky .................... Calibri Regular 10 pt, bezpatkové
%
% V tomto linuxovém prostředí nejsou Microsoft fonty dostupné, proto je
% výchozí volba metrická alternativa: Caladea ≈ Cambria, Carlito ≈ Calibri.
%
% Pokud kompilujete na systému s Microsoft fonty, pro nejvěrnější výstup
% nahraďte čtyři řádky níže těmito:
%   \setmainfont{Cambria}
%   \setsansfont{Calibri}
%   \newfontfamily\mgobodyfont{Cambria}
%   \newfontfamily\mgoheadingfont{Calibri}
%   \newfontfamily\mgotitlefont{Calibri}
%   \newfontfamily\mgofooterfont{Calibri}
% ---------------------------------------------------------

\setmainfont{Caladea}
\setsansfont{Carlito}
\newfontfamily\mgobodyfont{Caladea}
\newfontfamily\mgoheadingfont{Carlito}
\newfontfamily\mgotitlefont{Carlito}
\newfontfamily\mgofooterfont{Carlito}

\newcommand{\track}[2]{{\addfontfeatures{LetterSpace=#1}#2}}

% --- Rozměry stránky ---
\usepackage[
  a4paper,
  left=22mm,
  right=22mm,
  top=25mm,
  bottom=25mm,
  headheight=15pt,
  headsep=8mm,
  footskip=12mm
]{geometry}
\usepackage{tikz}
\usetikzlibrary{calc}
\usepackage[absolute,overlay]{textpos}

% --- Proměnné dokumentu ---
\newcommand{\typprace}{Seminární práce}
\newcommand{\nazevprace}{Název seminární práce}
\newcommand{\podnazevprace}{popř. podnázev seminární práce}
\newcommand{\autor}{jméno příjmení}
\newcommand{\skola}{Mendelovo gymnázium, Opava, příspěvková organizace}
\newcommand{\predmet}{Název předmětu}
\newcommand{\misto}{Opava}
\newcommand{\rok}{20\_\_}

% --- Odstavce: 1,5 řádkování, blok, mezera před odstavcem ---
\onehalfspacing
\setlength{\parindent}{12mm}
\setlength{\parskip}{9pt}
\AtBeginDocument{\justifying}

% První odstavec po nadpisu bez odsazení
\makeatletter
\patchcmd{\@afterheading}{\@afterindenttrue}{\@afterindentfalse}{}{}
\makeatother

% --- Záhlaví a zápatí ---
\usepackage{fancyhdr}
\usepackage{lastpage}
\pagestyle{fancy}
\fancyhf{}
\fancyhead[R]{{\mgobodyfont\itshape \typprace{} --- \nazevprace}}
\fancyfoot[C]{{\mgofooterfont\fontsize{10}{12}\selectfont \thepage/\pageref{LastPage}}}
\renewcommand{\headrulewidth}{0.4pt}
\renewcommand{\footrulewidth}{0.4pt}

\fancypagestyle{front}{%
  \fancyhf{}
  \renewcommand{\headrulewidth}{0pt}
  \renewcommand{\footrulewidth}{0pt}
}

% --- Nadpisy: bezpatkové, tučné, podobně jako Calibri Bold ---
\usepackage{titlesec}
\titleformat{\section}
  {\mgoheadingfont\fontsize{16}{19}\bfseries}
  {\thesection}{1.35em}{}
\titleformat{\subsection}
  {\mgoheadingfont\fontsize{14}{17}\bfseries}
  {\thesubsection}{1.2em}{}
\titleformat{\subsubsection}
  {\mgoheadingfont\fontsize{12}{14.5}\bfseries}
  {\thesubsubsection}{1.2em}{}

\titlespacing*{\section}{0pt}{18pt}{6pt}
\titlespacing*{\subsection}{0pt}{12pt}{3pt}
\titlespacing*{\subsubsection}{0pt}{9pt}{3pt}

% --- Obsah: tečky a odsazení pro dlouhé položky ---
\usepackage{tocloft}
\renewcommand{\contentsname}{Obsah}
\renewcommand{\cftsecleader}{\cftdotfill{\cftdotsep}}
\renewcommand{\cftsubsecleader}{\cftdotfill{\cftdotsep}}
\renewcommand{\cftsubsubsecleader}{\cftdotfill{\cftdotsep}}
\setlength{\cftsecnumwidth}{2.5em}
\setlength{\cftsubsecnumwidth}{3.2em}
\setlength{\cftsubsubsecnumwidth}{4.1em}
\setcounter{tocdepth}{3}
\setcounter{secnumdepth}{3}

% TOC v originálu: nadpis Calibri Bold 16 pt, položky Cambria 12 pt bez tučného řezu.
\renewcommand{\cfttoctitlefont}{\mgoheadingfont\fontsize{16}{19}\bfseries}
\renewcommand{\cftsecfont}{\normalfont}
\renewcommand{\cftsecpagefont}{\normalfont}
\renewcommand{\cftsubsecfont}{\normalfont}
\renewcommand{\cftsubsecpagefont}{\normalfont}
\renewcommand{\cftsubsubsecfont}{\normalfont}
\renewcommand{\cftsubsubsecpagefont}{\normalfont}


% --- Seznamy ---
\usepackage{enumitem}
\setlist[itemize]{label=--, leftmargin=12mm, itemsep=0pt, topsep=3pt}
\setlist[enumerate]{leftmargin=12mm, itemsep=0pt, topsep=3pt}

% --- Obrázky, tabulky, odkazy ---
\usepackage{graphicx}
\usepackage{booktabs}
\usepackage{array}
\usepackage{tabularx}
\usepackage[hidelinks]{hyperref}
\usepackage{url}
\usepackage{csquotes}
\newcolumntype{Y}{>{\raggedright\arraybackslash}X}

% --- Pomocné příkazy ---
\newcommand{\blankpage}{\newpage\thispagestyle{front}\null\newpage}

\begin{document}

% =========================================================
% Titulní strana - absolutní pozice podle původního PDF
% =========================================================
\thispagestyle{front}
\begin{titlepage}
\setlength{\TPHorizModule}{1mm}
\setlength{\TPVertModule}{1mm}

% Školní řádek: Cambria 16 pt Regular, patkové
\begin{textblock*}{210mm}(0mm,25mm)
\centering {\mgobodyfont\fontsize{16}{19}\selectfont\mdseries \skola}
\end{textblock*}

% Autor: Cambria 16 pt Regular, patkové, roztažené znaky
\begin{textblock*}{210mm}(0mm,80mm)
\centering {\mgobodyfont\fontsize{16}{19}\selectfont\mdseries\track{11}{\autor}}
\end{textblock*}

% Název: Calibri 28 pt Regular, bezpatkové, NE tučné, roztažené znaky
\begin{textblock*}{210mm}(0mm,115mm)
\centering {\mgotitlefont\fontsize{28}{33}\selectfont\mdseries\track{8}{\nazevprace}}
\end{textblock*}

% Podnázev: Calibri 16 pt Regular, bezpatkové, NE tučné
\begin{textblock*}{210mm}(0mm,129mm)
\centering {\mgotitlefont\fontsize{16}{19}\selectfont\mdseries\track{14}{\podnazevprace}}
\end{textblock*}

% Předmět: Cambria 16 pt Regular, patkové
\begin{textblock*}{210mm}(0mm,163mm)
\centering {\mgobodyfont\fontsize{16}{19}\selectfont\mdseries\track{6}{\typprace{} do předmětu \predmet}}
\end{textblock*}

% Místo a rok: Cambria 16 pt Regular, patkové
\begin{textblock*}{210mm}(0mm,256mm)
\centering {\mgobodyfont\fontsize{16}{19}\selectfont\mdseries\track{12}{\misto{} \rok}}
\end{textblock*}

\end{titlepage}

\blankpage

% =========================================================
% Prohlášení
% =========================================================
\thispagestyle{front}
\vspace*{12cm}
Prohlašuji, že jsem tuto seminární práci vypracoval(a) samostatně na základě uvedených pramenů a literatury.

\vspace{1.5cm}
\noindent V~\misto{} dne \today \hfill Podpis: \rule{45mm}{0.4pt}

\blankpage

% =========================================================
% Obsah
% =========================================================
\thispagestyle{front}
\tableofcontents
\newpage

% =========================================================
% Vlastní text práce
% =========================================================
\section{Nadpis první úrovně}
První odstavec za nadpisem začíná bez odsazení. Text je zarovnán do bloku, má velikost 12~pt, řádkování 1,5 a odstavce mají menší mezeru před sebou. Sem vložte úvodní část seminární práce.

Další odstavce už mají odsazený první řádek. Tato šablona je určena jako jednoduchý základ pro seminární práci: titulní strana, prohlášení, obsah, číslované nadpisy, záhlaví, zápatí, poznámky pod čarou, seznamy, obrázky, tabulky a seznam použitých zdrojů.

\section{Dlouhý nadpis první úrovně, který přesahuje délku jednoho řádku a ukazuje, jak se nadpis zalomí}
Ukázkový text kapitoly. Vlastní obsah nahraďte svým textem.

\subsection{Nadpis druhé úrovně}
Text podkapitoly.

\subsection{Dlouhý nadpis druhé úrovně, který přesahuje délku jednoho řádku a ukazuje zalomení}
Text další podkapitoly.

\subsubsection{Nadpis třetí úrovně}
Text části třetí úrovně.

\section{Seznamy}
Ukázka nečíslovaného seznamu:
\begin{itemize}
  \item první položka,
  \item druhá položka,
  \item třetí položka.
\end{itemize}

Ukázka víceúrovňového číslovaného seznamu:
\begin{enumerate}[label=\arabic*)]
  \item První položka
  \begin{enumerate}[label=\alph*)]
    \item Vnořená položka
    \item Vnořená položka
    \begin{enumerate}[label=\roman*)]
      \item Další úroveň
      \item Další úroveň
    \end{enumerate}
  \end{enumerate}
  \item Druhá položka
\end{enumerate}

\section{Obrázky}
Obrázek lze vložit standardním prostředím \texttt{figure}. Popisek je pod obrázkem a na obrázek se lze odkazovat v textu, například Obrázek~\ref{fig:ukazka}.

\begin{figure}[h]
  \centering
  \fbox{\rule{0pt}{45mm}\rule{0.75\textwidth}{0pt}}
  \caption{Ukázkový prostor pro obrázek.}
  \label{fig:ukazka}
\end{figure}

\section{Tabulky}
Tabulka může mít jednoduchý akademický vzhled s výraznější hlavičkou.

\begin{table}[h]
  \centering
  \caption{Ukázková tabulka}
  \begin{tabularx}{\textwidth}{|Y|Y|Y|Y|}
    \hline
    \textbf{1. buňka} & \textbf{2. buňka} & \textbf{3. buňka} & \textbf{4. buňka} \\
    \hline
    Text & Text & Text & Text \\
    \hline
    Text & Text & Text & Text \\
    \hline
  \end{tabularx}
  \label{tab:ukazka}
\end{table}

\section{Citace a poznámky}
Krátkou poznámku lze vložit pod čarou.\footnote{Toto je ukázková poznámka pod čarou.} Odkaz na zdroj lze psát ručně, například podle autora a roku vydání, nebo využít citační balíčky podle požadavků školy.

\section{Použité zdroje}
\begin{thebibliography}{9}
  \bibitem{zdroj1}
  PŘÍJMENÍ, Jméno. \textit{Název knihy}. Místo vydání: Nakladatel, rok. ISBN.
\end{thebibliography}

\end{document}
